\documentclass[../../main.tex]{subfiles}% 注意这里的文件路径不能用 ./main.tex ,否则用latexmk编译子文件会报错
\graphicspath{{\subfix{./image/}}{\subfix{../../image/}}} % 指定图片引用路径,后续可以直接使用图片文件名
% 如果图片存放在上级目录,则使用 ../../image/ 作为备用路径

% 例如:
% \begin{figure}[H]
% \centering
% \includegraphics[width=0.3\textwidth]{图.png}
% \caption{}
% \label{figure:图}
% \end{figure}
% 注意:上述\label{}一定要放在\caption{}之后,否则引用图片序号会只会显示??.

\begin{document}

\section{互素多项式的应用}

\begin{proposition}
设 $f(x), g(x)$ 是数域 $\mathbb{K}$ 上的互素多项式，$A$ 是 $\mathbb{K}$ 上的 $n$ 阶方阵，满足 $f(A) = O$，证明：$g(A)$ 是可逆矩阵。
\end{proposition}
\begin{proof}
根据假设，存在 $\mathbb{K}$ 上的多项式 $u(x), v(x)$，使得
\begin{align*}
f(x)u(x) + g(x)v(x) = 1。
\end{align*}
在上述式中代入 $x = A$，可得恒等式
\begin{align*}
f(A)u(A) + g(A)v(A) = I_n。
\end{align*}
因为 $f(A) = O$，故有 $g(A)v(A) = I_n$，从而 $g(A)$ 是非零矩阵且 $g(A)^{-1} = v(A)$.
\end{proof}

\begin{proposition}
设 $f(x), g(x)$ 是数域 $\mathbb{K}$ 上的互素多项式，$A$ 是 $\mathbb{K}$ 上的 $n$ 阶方阵，证明：$f(A)g(A) = O$ 的充要条件是 $r(f(A)) + r(g(A)) = n$。
\end{proposition}
\begin{proof}
根据假设，存在 $\mathbb{K}$ 上的多项式 $u(x), v(x)$，使得
\begin{align*}
f(x)u(x) + g(x)v(x) = 1。
\end{align*}
在上述式中代入 $x = A$，可得恒等式
\begin{align*}
f(A)u(A) + g(A)v(A) = I_n。
\end{align*}
考虑如下分块矩阵的初等变换：
\begin{align*}
\begin{pmatrix}
f(A) & O \\
O & g(A)
\end{pmatrix}
&\to
\begin{pmatrix}
f(A) & f(A)u(A) \\
O & g(A)
\end{pmatrix}
\to
\begin{pmatrix}
f(A) & I_n \\
O & g(A)
\end{pmatrix}
\to
\begin{pmatrix}
f(A) & I_n \\
-f(A)g(A) & O
\end{pmatrix}
\to
\begin{pmatrix}
O & I_n \\
-f(A)g(A) & O
\end{pmatrix},
\end{align*}
故有 $r(f(A)) + r(g(A)) = r(f(A)g(A)) + n$，从而结论得证。
\end{proof}

\begin{proposition}\label{proposition:互素多项式诱导直和分解1}
设 $f(x), g(x)$ 是数域 $\mathbb{K}$ 上的互素多项式，$\varphi$ 是 $\mathbb{K}$ 上 $n$ 维线性空间 $V$ 上的线性变换，满足 $f(\varphi)g(\varphi) = 0$，证明：$V = V_1 \oplus V_2$，其中 $V_1 = \text{Ker} f(\varphi), V_2 = \text{Ker} g(\varphi)$。
\end{proposition}
\begin{remark}
这个命题告诉我们:多项式的互素因式分解可以诱导出空间的直和分解,从几何层面上看,这就是相似标准型理论原始的出发点.
\end{remark}
\begin{proof}
根据假设，存在 $\mathbb{K}$ 上的多项式 $u(x), v(x)$，使得
\begin{align*}
f(x)u(x) + g(x)v(x) = 1。
\end{align*}
在上述式中代入 $x = \varphi$，可得恒等式
\begin{align*}
f(\varphi)u(\varphi) + g(\varphi)v(\varphi) = I_V。
\end{align*}
对任意的 $\alpha \in V$，由上述可得
\begin{align*}
\alpha = f(\varphi)u(\varphi)(\alpha) + g(\varphi)v(\varphi)(\alpha),
\end{align*}
注意到 
\begin{align*}
g\left( \varphi \right) \left( f\left( \varphi \right) u\left( \varphi \right) \left( \alpha \right) \right) =g\left( \varphi \right) f\left( \varphi \right) u\left( \varphi \right) \left( \alpha \right) =u\left( \varphi \right) f\left( \varphi \right) g\left( \varphi \right) \left( \alpha \right) =0,
\\
f\left( \varphi \right) \left( g\left( \varphi \right) u\left( \varphi \right) \left( \alpha \right) \right) =f\left( \varphi \right) g\left( \varphi \right) u\left( \varphi \right) \left( \alpha \right) =u\left( \varphi \right) f\left( \varphi \right) g\left( \varphi \right) \left( \alpha \right) =0.
\end{align*}
于是$f(\varphi)u(\varphi)(\alpha) \in \text{Ker} g(\varphi), g(\varphi)v(\varphi)(\alpha) \in \text{Ker} f(\varphi)$，故有 $V = V_1 + V_2$。任取 $\beta \in V_1 \cap V_2$，由上述可得
\begin{align*}
\beta = u(\varphi)f(\varphi)(\beta) + v(\varphi)g(\varphi)(\beta) = 0,
\end{align*}
故有 $V_1 \cap V_2 = 0$，因此 $V = V_1 \oplus V_2$。     
\end{proof}

\begin{theorem}\label{theorem:线性变换互素分解对应核子空间直和分解}
设 $V$ 是域 $\mathbb{F}$ 上的线性空间，$A$ 是 $V$ 上线性变换，$f, f_i \in \mathbb{F}[x], i = 1,2,\cdots, s$ 满足 $f = f_1f_2\cdots f_s$ 并有 $f_1, f_2, \cdots, f_s$ 两两互素. 则
$$\mathrm{Ker} f(A) = \mathrm{Ker} f_1(A) \oplus \mathrm{Ker} f_2(A) \oplus \cdots \oplus \mathrm{Ker} f_s(A).$$
\end{theorem}
\begin{remark}
这个定理最常见的情况是:当$f$为$A$的零化多项式时,就有$\mathrm{Ker} f(A)=V$,此时利用这个定理就能得到一个全空间的直和分解.
\end{remark}
\begin{proof}
当 $s = 2$, 此时由裴蜀等式得 $p_1, p_2 \in \mathbb{F}[x]$ 使得 $p_1f_1 + p_2f_2 = 1$. 于是对 $\alpha \in \mathrm{Ker} f(A)$, 我们有
$$\alpha = p_1(A)f_1(A)\alpha + p_2(A)f_2(A)\alpha.$$
定义
$$\alpha_1 \triangleq p_2(A)f_2(A)\alpha,\quad \alpha_2 \triangleq p_1(A)f_1(A)\alpha.$$
现在
$$f_1(A)\alpha_1 = p_2(A)f_1(A)f_2(A)\alpha = p_2(A)f(A)\alpha = 0 \Rightarrow \alpha_1 \in \mathrm{Ker} f_1(A);$$
$$f_2(A)\alpha_2 = p_1(A)f_1(A)f_2(A)\alpha = p_1(A)f(A)\alpha = 0 \Rightarrow \alpha_2 \in \mathrm{Ker} f_2(A).$$
因此我们的确有
$$\mathrm{Ker} f(A) = \mathrm{Ker} f_1(A) + \mathrm{Ker} f_2(A).$$
现在设 $\alpha \in \mathrm{Ker} f_1(A) \cap \mathrm{Ker} f_2(A)$, 于是 $\alpha = p_1(A)f_1(A)\alpha + p_2(A)f_2(A)\alpha = 0$, 即
$$\mathrm{Ker} f(A) = \mathrm{Ker} f_1(A) \oplus \mathrm{Ker} f_2(A).$$
对 $s > 2$, 我们考虑归纳法. 假设命题对 $s - 1$ 已经成立, 设 $g(x) \triangleq f_2(x)f_3(x)\cdots f_s(x)$, 则由\refpro{proposition:两两互素的多项式组的乘积也互素}知$f_1, g$ 互素. 由 $s = 2$ 的结论知
$$\mathrm{Ker} f(A) = \mathrm{Ker} f_1(A) \oplus \mathrm{Ker} g(A).$$
对 $g$ 用归纳假设得
$$\mathrm{Ker} f(A) = \mathrm{Ker} f_1(A) \oplus \mathrm{Ker} f_2(A) \oplus \cdots \oplus \mathrm{Ker} f_s(A).$$
\end{proof}

\begin{proposition}\label{proposition:互素分式多项式诱导的值域直和分解}
设 $V$ 是域 $\mathbb{F}$ 上的线性空间, $A$ 是 $V$ 上线性变换. 设 $f, f_i \in \mathbb{F}[x], i = 1, 2, \cdots, s$ 满足 $f = f_1f_2\cdots f_s$ 并有 $f_1, f_2, \cdots, f_s$ 两两互素, $f$ 是 $A$ 的零化多项式. 证明:
$$
V = \mathrm{Im} \frac{f(A)}{f_1(A)} \oplus \mathrm{Im} \frac{f(A)}{f_2(A)} \oplus \cdots \oplus \mathrm{Im} \frac{f(A)}{f_s(A)}.
$$
\end{proposition}
\begin{proof}
由\refthe{theorem:线性变换互素分解对应核子空间直和分解},我们有
$$
\mathrm{Ker} f(A) = \mathrm{Ker} f_1(A) \oplus \mathrm{Ker} f_2(A) \oplus \cdots \oplus \mathrm{Ker} f_s(A).
$$
设
$$
\alpha = \frac{f(A)}{f_i(A)} \beta \in \mathrm{Im} \frac{f(A)}{f_i(A)},
$$
我们有
$$
f_i(A) \alpha = f(A) \beta = 0 \implies \alpha \in \mathrm{Ker} f_i(A),
$$
故$\mathrm{Im}\frac{f\left( A \right)}{f_i\left( A \right)}\subseteq \mathrm{ker}f_i\left( A \right) ,i=1,2,\cdots ,s.$
于是由\refpro{proposition:直和分解的每个子空间的和还是直和}知
$$
\mathrm{Im} \frac{f(A)}{f_1(A)} + \mathrm{Im} \frac{f(A)}{f_2(A)} + \cdots + \mathrm{Im} \frac{f(A)}{f_s(A)}
$$
是直和.

由\refpro{proposition:f除以两两互素多项式得到的多项式组互素}知 $\frac{f}{f_i}, i = 1, 2, \cdots, s$ 互素, 由裴蜀等式我们知道存在 $u_i \in \mathbb{F}[x], i = 1, 2, \cdots, s$ 使得
$$
\sum_{i=1}^s u_i \frac{f}{f_i} = 1.
$$
于是对 $\alpha \in V$, 我们有
$$
\alpha = \sum_{i=1}^s \frac{f(A)}{f_i(A)} (u_i(A) \alpha) \in \mathrm{Im} \frac{f(A)}{f_1(A)} \oplus \mathrm{Im} \frac{f(A)}{f_2(A)} \oplus \cdots \oplus \mathrm{Im} \frac{f(A)}{f_s(A)}.
$$
现在我们知道
$$
V = \mathrm{Im} \frac{f(A)}{f_1(A)} \oplus \mathrm{Im} \frac{f(A)}{f_2(A)} \oplus \cdots \oplus \mathrm{Im} \frac{f(A)}{f_s(A)}.
$$
\end{proof}

\begin{corollary}
设$V$是数域$\mathrm{P}$上的线性空间，$A$是空间$V$上的线性变换。$f(x),$$f_1(x)$, $\cdots,$$f_s(x)$ $\in $ $\mathrm{P}[x],$ $f(x)$ $=$ $f_1(x)f_2(x)\cdots f_s(x)$是$A$的零化多项式，且$f_1(x), \cdots, f_s(x)$两两互素，现在设$F_i(x) = \frac{f(x)}{f_i(x)}$$(i = 1, 2, \cdots, s)$，则对于任意的$i = 1, 2, \cdots, s$有
\begin{align*}
\text{Ker} f_i(A) = \text{Im}\, F_i(A).
\end{align*}
\end{corollary}
\begin{proof}
任取$\alpha \in \text{Ker}f_i(A)$,由\refpro{proposition:两两互素的多项式组的乘积也互素}知$f_i$与$F_i$互素,故存在$u_i,v_i\in \mathrm{P}[x]$,使得
\begin{align*}
u_i(x)f_i(x)+v_i(x)F_i(x)=1.
\end{align*}
从而
\begin{align*}
\alpha=u_i(A)f_i(A)(\alpha) + v_i(A)F_i(A)(\alpha) =  F_i(A)(v_i(A)(\alpha)).
\end{align*}
故$\alpha \in \text{Im}F_i(x)$,因此
\begin{align}\label{eq::9-k3434g2g}
\text{Ker}f_i(A)\subset \text{Im}F_i(A),i=1,2,\cdots,s.
\end{align}
再任取$F_i(A)\alpha \in \text{Im}\, F_i(A),$则$f_i(A)((F_i(A))\alpha)=f(A)\alpha=0,$因此$F_i(A)\alpha\in \text{Ker}f_i(A)$,故$\text{Im}F_i(A)\subseteq \text{Ker}f_i(A)$.
综上可知
\begin{align*}
\text{Ker} f_i(A) = \text{Im}\, F_i(A),\,i=1,2,\cdots,s.
\end{align*}
\end{proof}

\begin{example}
设 $B_i \in \mathbb{C}^{n \times n}, i = 1, 2, \cdots, k$ 是幂等矩阵, $A = B_1 B_2 \cdots B_k$, 证明:
$$
r(I - A) \leqslant k(n - r(A)).
$$
\end{example}
\begin{proof}
将$A$,$B_i$都看作在$\mathbb{C}^n$上左乘诱导的线性变换.由$B_i$是幂等矩阵知
\begin{align*}
B^2 = B \Longleftrightarrow \mathrm{Ker} (B_i^2 - B_i) = \mathbb{C}^n,\ i = 1,2,\cdots,k.
\end{align*}
由\refthe{theorem:线性变换互素分解对应核子空间直和分解}知
\begin{align*}
\mathbb{C}^n = \mathrm{Ker} (B_i^2 - B_i) = \mathrm{Ker} B_i \oplus \mathrm{Ker} (I - B_i),\ i = 1,2,\cdots,k.
\end{align*}
再结合\hyperref[theorem:线性映射的维数公式]{维数公式}可得
\begin{align*}
n &= \dim \mathrm{Ker} B_i + \dim \mathrm{Ker} (I - B_i) = n - \dim \mathrm{Im} B_i + n - \dim \mathrm{Im} (I - B_i) \\
&= n - r(B_i) + n - r(I - B_i),\ i = 1,2,\cdots,k.
\end{align*}
因此
\begin{align*}
r(I - B_i) + r(B_i) = n,\ i = 1,2,\cdots,k. 
\end{align*}
注意到
\begin{align*}
I - A = I - B_1B_2\cdots B_k = (I - B_1)B_2\cdots B_k + (I - B_2)B_3\cdots B_k + \cdots + (I - B_{k-1})B_k + I - B_k,
\end{align*}
故由\hyperref[proposition:矩阵秩的基本公式]{矩阵秩的基本公式-(2)}知
\begin{align*}
r(A) = r(B_1B_2\cdots B_k) \leqslant r(B_i),\ i = 1,2,\cdots,k.
\end{align*}
从而再由\hyperref[proposition:矩阵秩的基本公式]{矩阵秩的基本公式-(2)和(5)}知
\begin{align*}
r(I - A) &\leqslant r(I - B_1) + r(I - B_2) + \cdots + r(I - B_k) \\
&= kn - r(B_1) - r(B_2) - \cdots - r(B_k) \\
&\leqslant kn - kr(A) = k(n - r(A)).
\end{align*}
\end{proof}

\begin{example}
设 $\mathbb{Q}(\sqrt[n]{2}) = \{a_0 + a_1 \sqrt[n]{2} + a_2 \sqrt[4]{4} + \cdots + a_{n-1} \sqrt[n]{2^{n-1}} \mid a_i \in \mathbb{Q}, 0 \leqslant  i \leqslant  n - 1\}$，证明：$\mathbb{Q}(\sqrt[n]{2})$ 是一个数域，并求 $\mathbb{Q}(\sqrt[n]{2})$ 作为 $\mathbb{Q}$ 上线性空间的一组基。
\end{example}
\begin{proof}
设 $f(x) = x^n - 2$，由 Eisenstein 判别法可知 $f(x)$ 在 $\mathbb{Q}$ 上不可约，从而 $f(x)$ 是 $\sqrt[n]{2}$ 的极小多项式。我们先证明：$a_0 + a_1 \sqrt[n]{2} + \cdots + a_{n-1} \sqrt[n]{2^{n-1}} = 0$ 的充要条件是 $a_0 = a_1 = \cdots = a_{n-1} = 0$。充分性是显然的，现证必要性：令 $g(x) = a_0 + a_1 x + \cdots + a_{n-1} x^{n-1}$，则 $g(\sqrt[n]{2}) = 0$，由极小多项式的基本性质可得 $f(x) \mid g(x)$。因为 $g(x)$ 的次数小于 $n$，故只能是 $g(x) = 0$，即 $a_0 = a_1 = \cdots = a_{n-1} = 0$。

利用 $\sqrt[n]{2} = 2$ 容易验证，$\mathbb{Q}(\sqrt[n]{2})$ 中任意两个数的加法、减法和乘法都是封闭的。要证明 $\mathbb{Q}(\sqrt[n]{2})$ 是数域，只要证明除法或者取倒数封闭即可。任取 $\mathbb{Q}(\sqrt[n]{2})$ 中的非零数 $\alpha = a_0 + a_1 \sqrt[2]{2} + \cdots + a_{n-1} \sqrt[n]{2^{n-1}} \neq 0$，由上面的讨论可知 $a_0, a_1, \cdots, a_{n-1}$ 不全为零。令 $g(x) = a_0 + a_1 x + \cdots + a_{n-1} x^{n-1}$，则 $g(\sqrt[n]{2}) \neq 0$。因为 $f(x)$ 不可约且 $g(x) \neq 0$ 的次数小于 $n$，故 $f(x)$ 与 $g(x)$ 互素，由\hyperref[theorem:多项式互素的充要条件]{多项式互素的充要条件}可知，存在有理系数多项式 $u(x), v(x)$，使得
\begin{align*}
f(x)u(x) + g(x)v(x) = 1,
\end{align*}
在上述中代入 $x = \sqrt[n]{2}$，可得 $\sqrt[n]{2} v(\sqrt[n]{2}) = 1$，于是 $\alpha^{-1} = v(\sqrt[n]{2}) \in \mathbb{Q}(\sqrt[n]{2})$。因此，$\mathbb{Q}(\sqrt[n]{2})$ 是数域。

由 $\mathbb{Q}(\sqrt[n]{2})$ 的定义可知，$\mathbb{Q}(\sqrt[n]{2})$ 中任一元都是 $1, \sqrt[2]{2}, \cdots, \sqrt[n]{2^{n-1}}$ 的 $\mathbb{Q}$-线性组合；又由开始的讨论可知，$1, \sqrt[2]{2}, \cdots, \sqrt[n]{2^{n-1}}$ 是 $\mathbb{Q}$-线性无关的，因此它们构成了 $\mathbb{Q}(\sqrt[n]{2})$ 作为 $\mathbb{Q}$ 上线性空间的一组基。特别地，$\dim_{\mathbb{Q}} \mathbb{Q}(\sqrt[n]{2}) = n$。
\end{proof}

\begin{proposition}\label{proposition:不可约多项式与向量组线性无关}
\begin{enumerate}[(1)]
\item 设 $f(x) = x^n + a_1 x^{n-1} + \cdots + a_{n-1} x + a_n$ 是数域 $\mathbb{K}$ 上的不可约多项式，$\varphi$ 是 $\mathbb{K}$ 上 $m(\geqslant n)$ 维线性空间 $V$ 上的线性变换，$\alpha_1 \neq 0, \alpha_2, \cdots, \alpha_n$ 是 $V$ 中的向量，满足
\begin{align*}
\varphi (\alpha _1)=\alpha _2,\varphi (\alpha _2)=\alpha _3,\cdots ,\varphi (\alpha _{n-1})=\alpha _n,\varphi (\alpha _n)=-a_n\alpha _1-a_{n-1}\alpha _2-\cdots -a_1\alpha _n.
\end{align*}
证明：$\{\alpha_1, \alpha_2, \cdots, \alpha_n\}$ 线性无关.

\item 设 $f(x) = x^n + a_1 x^{n-1} + \cdots + a_{n-1} x + a_n$ 是数域 $\mathbb{K}$ 上的不可约多项式, $\varphi$ 是 $\mathbb{K}$ 上线性空间 $V$ 的一个线性变换。若存在非零向量 $\alpha \in V$ 使得
$f(\varphi)(\alpha) = 0,$
即
\begin{align*}
\varphi^n(\alpha) + a_1 \varphi^{n-1}(\alpha) + \cdots + a_{n-1}\varphi(\alpha) + a_n \alpha = 0,
\end{align*}
则向量组$\{\alpha, \varphi(\alpha), \varphi^2(\alpha), \cdots, \varphi^{n-1}(\alpha)\}$
必线性无关.
\end{enumerate}
\end{proposition}
\begin{proof}
(1)和(2)等价,只证明(1).用反证法，设存在不全为零的 $n$ 个数 $c_1, c_2, \cdots, c_n$，使得
\begin{align*}
c_1 \alpha_1 + c_2 \alpha_2 + \cdots + c_n \alpha_n = 0,
\end{align*}
则有
\begin{align*}
(c_1 I_V + c_2 \varphi + \cdots + c_n \varphi^{n-1})(\alpha_1) = 0。
\end{align*}
令 $g(x) = c_1 + c_2 x + \cdots + c_n x^{n-1}$，则 $g(x) \neq 0$ 且 $g(\varphi)(\alpha_1) = 0$。另一方面，由条件容易验证 $f(\varphi)(\alpha_1) = 0$。因为 $f(x)$ 不可约且 $g(x)$ 的次数小于 $n$，故 $f(x)$ 与 $g(x)$ 互素，从而存在 $\mathbb{K}$ 上的多项式 $u(x), v(x)$，使得
\begin{align*}
f(x)u(x) + g(x)v(x) = 1。
\end{align*}
在上述中代入 $x = \varphi$，可得恒等式
\begin{align*}
f(\varphi)u(\varphi) + g(\varphi)v(\varphi) = I_V。
\end{align*}
上式两边同时作用 $\alpha_1$ 可得
\begin{align*}
\alpha_1 = u(\varphi)f(\varphi)(\alpha_1) + v(\varphi)g(\varphi)(\alpha_1) = 0,
\end{align*}
这与$\alpha_1\ne 0$矛盾!
\end{proof}

\begin{example}\label{example:有理数域上线性变换与向量组线性无关的命题}
设$\alpha, \beta, \gamma \in \mathbb{Q}^3$且$\alpha \neq 0$. 若$A \in \mathbb{Q}^{3 \times 3}$满足
\begin{align*}
A\alpha =\beta ,A\beta =-\gamma ,A\gamma =\alpha -\beta .
\end{align*}
证明:$\alpha, \beta, \gamma$在$\mathbb{Q}$上线性无关.
\end{example}
\begin{proof}

{\color{blue}证法一:}取$\alpha'=-\alpha,\beta'=-\beta,\gamma'=\gamma$,则
\begin{align*}
A\alpha =\beta ,A\beta =-\gamma ,A\gamma =\alpha -\beta \Longleftrightarrow A\alpha' =\beta' ,A\beta' =\gamma' ,A\gamma' =\alpha' +\beta' .
\end{align*}
由\refpro{proposition:不可约多项式与向量组线性无关}可知$\alpha',\beta',\gamma'$线性无关,而$\{\alpha,\beta,\gamma\}$和$\{\alpha',\beta',\gamma'\}$等价,故$\alpha,\beta,\gamma$也线性无关.

{\color{blue}证法二:}
若$\alpha, \beta$在$\mathbb{Q}$上线性相关，则存在$k\in \mathbb{Q}$，使得$\beta =k\alpha$。从而由条件可得
\begin{align*}
A\alpha =\beta =k\alpha \Longrightarrow A\beta =kA\alpha =k^2\alpha =-\gamma 
\Longrightarrow A\gamma =\alpha -\beta =\left( 1-k \right) \alpha =-k^2A\alpha =-k^3\alpha .
\end{align*}
于是就有$\left( k^3-k+1 \right) \alpha =0$。又$\alpha \ne 0$，故$k^3-k+1=0$。但这个方程没有有理根，矛盾！故$\alpha, \beta$在$\mathbb{Q}$上线性无关。

若$\alpha, \beta, \gamma$在$\mathbb{Q}$上线性相关，则存在$a,b\in \mathbb{Q}$，使得$\gamma =a\alpha +b\beta$。由条件可得
\begin{align*}
A\gamma =\alpha -\beta =aA\alpha +bA\beta =a\beta -b\gamma =a\beta -b\left( a\alpha +b\beta \right) =-ab\alpha +\left( a+b^2 \right) \beta .
\end{align*}
因此
\begin{align*}
ab=1,\quad a+b^2=-1\Longrightarrow a+\frac{1}{a^2}=-1\Longrightarrow a^3-a^2+1=0,
\end{align*}
而上式无有理根，矛盾！故$\alpha, \beta, \gamma$在$\mathbb{Q}$上线性无关。

{\color{blue}证法三:}
记$A$左乘诱导的线性变换为$\mathscr{A}$,矩阵 $\Lambda = \begin{pmatrix} 0 & 0 & 1 \\ 1 & 0 & -1 \\ 0 & -1 & 0 \end{pmatrix}$，那么
\begin{align*}
\mathscr{A}(\alpha, \beta, \gamma) = (\alpha, \beta, \gamma)\Lambda.
\end{align*}
若 $\alpha, \beta, \gamma$ 线性相关，则存在非零向量 $x \in \mathbb{Q}^3$ 使得 $(\alpha, \beta, \gamma)x = 0$，反复作用 $\mathscr{A}$ 可得
\begin{align*}
(\alpha, \beta, \gamma)x = 0, \quad (\alpha, \beta, \gamma)\Lambda x = 0, \quad (\alpha, \beta, \gamma)\Lambda ^2x = 0.
\end{align*}
若 $x, \Lambda x, \Lambda^2x$ 线性相关，则存在不全为零的有理数 $a, b, c$，使得
\begin{align*}
ax + b\Lambda x + c\Lambda ^2x = (aE + b\Lambda  + c\Lambda ^2)x = 0.
\end{align*}
记 $g(x) = a+bx+cx^2$，上式说明 $0$ 是 $g(\Lambda)$ 的特征值，因此 $g(\Lambda)$ 不可逆。而实际上 $\Lambda$ 的特征多项式为 $f(x) = x^3 - x + 1$，其没有有理根，因此在有理数域上不可约，同时 $f(x) \nmid g(x)$，所以 $(f(x), g(x)) = 1$，那么存在 $u(x), v(x) \in \mathbb{Q}[x]$，使得
\begin{align*}
u(x)f(x) + v(x)g(x) = 1,
\end{align*}
进而
\begin{align*}
u(\Lambda )f(\Lambda ) + v(\Lambda )g(\Lambda ) = v(\Lambda )g(\Lambda ) = E,
\end{align*}
即 $g(\Lambda)$ 可逆，矛盾!于是 $x, \Lambda x, \Lambda^2x$ 线性无关，进而它们构成 $\mathbb{Q}^3$ 的一组基，那么存在 $a_0, b_0, c_0 \in \mathbb{Q}$，使得 $a_0x + b_0\Lambda x + c_0\Lambda ^2x = (1, 0, 0)^T$，进而
\begin{align*}
\alpha = (\alpha, \beta, \gamma)(1, 0, 0)^T = (\alpha, \beta, \gamma)(a_0x + b_0\Lambda x + c_0\Lambda ^2x) = 0.
\end{align*}
这与已知矛盾!因此 $\alpha, \beta, \gamma$ 在 $\mathbb{Q}$ 上线性无关.
\end{proof}

\begin{example}
设$V$是$\mathbb{Q}$上4维空间，$\varphi$是$V$上的线性变换. 若
$$\alpha_i \in V, i = 1,2,3,4,5.$$
且满足
\begin{gather*}
\alpha _1\ne 0,\quad \alpha _4\ne \alpha _1+\alpha _2,\quad \varphi( \alpha _1)=\alpha _2,\quad \varphi (\alpha _2)=\alpha _3,
\\
\varphi (\alpha _3)=\alpha _1+\alpha _2,\quad \varphi (\alpha _4)=\alpha _5,\quad \varphi (\alpha _5)=\alpha _3+\alpha _4.
\end{gather*}
求$\det \varphi$.
\end{example}
\begin{proof}
由\refpro{proposition:不可约多项式与向量组线性无关}可知
$\alpha_1,\alpha_2,\alpha_3$线性无关.若$\alpha_4\in \mathrm{span}\{\alpha_1,\alpha_2,\alpha_3\}$,设
$$\alpha_4=a\alpha_1+b\alpha_2+c\alpha_3,\quad a,b,c\in \mathbb{Q}.$$
由条件可得
\begin{align*}
\varphi (\alpha_4)=\alpha_5\Longrightarrow \varphi^2(\alpha_4)=\varphi (\alpha_5)=\alpha_3+\alpha_4=a\alpha_1+b\alpha_2+(c+1)\alpha_3,
\end{align*}
\begin{align*}
\varphi^2(\alpha_4)=\varphi^2(a\alpha_1+b\alpha_2+c\alpha_3)=\varphi\left(c\alpha_1+(a+c)\alpha_2+b\alpha_3\right)=b\alpha_1+(b+c)\alpha_2+(a+c)\alpha_3.
\end{align*}
从而
\begin{align*}
a=b,\quad b=b+c,\quad c+1=a+c\Longrightarrow a=b=1,\quad c=0.
\end{align*}
故$\alpha_4=\alpha_1+\alpha_2$与条件矛盾!因此$\alpha_1,\alpha_2,\alpha_3,\alpha_4$线性无关.又因为$V$是$\mathbb{Q}$上$4$维空间,所以$\alpha_1,\alpha_2,\alpha_3,\alpha_4$就是$V$的一组基.从而
\begin{align*}
\varphi\left( \alpha_1,\alpha_2,\alpha_3,\alpha_4 \right)=\left( \alpha_1,\alpha_2,\alpha_3,\alpha_4 \right)\begin{pmatrix}
0&		0&		1&		x\\
1&		0&		1&		y\\
0&		1&		0&		z\\
0&		0&		0&		m\\
\end{pmatrix},
\end{align*}
其中$x,y,z,m\in \mathbb{Q}$.由条件可知
\begin{align*}
\varphi^2(\alpha_4)=\varphi (\alpha_5)=\alpha_3+\alpha_4.
\end{align*}
于是$\varphi^2(\alpha_4)$的在基$\alpha_1,\alpha_2,\alpha_3,\alpha_4$下的坐标就是
\begin{align*}
\begin{pmatrix}
0&		0&		1&		x\\
1&		0&		1&		y\\
0&		1&		0&		z\\
0&		0&		0&		m\\
\end{pmatrix}^2 \left( \begin{array}{c}
0\\
0\\
0\\
1\\
\end{array} \right) =\left( \begin{array}{c}
0\\
0\\
1\\
1\\
\end{array} \right) \Longleftrightarrow \left( \begin{array}{c}
z+xm\\
x+z+ym\\
y+zm\\
m^2\\
\end{array} \right) =\left( \begin{array}{c}
0\\
0\\
1\\
1\\
\end{array} \right) .
\end{align*}
解得$(x,y,z,m)^T=(-1,0,1,1)^T$或$(1,2,1,-1)^T$.故
\begin{align*}
\det \varphi=\begin{vmatrix}
0&		0&		1&		-1\\
1&		0&		1&		0\\
0&		1&		0&		1\\
0&		0&		0&		1\\
\end{vmatrix}=-1\text{或}\begin{vmatrix}
0&		0&		1&		1\\
1&		0&		1&		2\\
0&		1&		0&		1\\
0&		0&		0&		-1\\
\end{vmatrix}=1.
\end{align*}
\end{proof}

\begin{example}
设 $V$ 是有理数域 $\mathbb Q$ 上的 $n$ 维线性空间，证明：存在 $V$ 上 $n$ 个线性变换
$\varphi_1,\varphi_2,\cdots,\varphi_n$，使得对任意非零向量 $\alpha\in V$，向量组
\[
\{\varphi_1(\alpha),\varphi_2(\alpha),\cdots,\varphi_n(\alpha)\}
\]
均线性无关。
\end{example}
\begin{proof}

取有理系数多项式
\[
f(x)=x^n-2\in\mathbb Q[x].
\]
由\hyperref[theorem:Eisenstein判别法]{Eisenstein判别法}知$f(x)$ 在 $\mathbb Q$ 上不可约。

取 $V$ 的一组基 $\{e_1,e_2,\cdots,e_n\}$，在 $V$ 上定义线性变换 $\varphi$ 为
\[
\varphi(e_1)=e_2,\quad \varphi(e_2)=e_3,\quad \cdots,\quad \varphi(e_{n-1})=e_n,\quad \varphi(e_n)=2e_1.
\]
直接计算可得
\[
\varphi^n(e_i)=2e_i\quad (i=1,2,\cdots,n),
\]
故 $\varphi^n=2I$，即
\[
f(\varphi)=\varphi^n-2I=0.
\]
因此对任意向量 $\alpha\in V$，都有 $f(\varphi)(\alpha)=0$。

现在定义
\[
\varphi_1=I,\quad \varphi_2=\varphi,\quad \varphi_3=\varphi^2,\quad \cdots,\quad \varphi_n=\varphi^{n-1}.
\]
显然 $\varphi_1,\varphi_2,\cdots,\varphi_n$ 都是 $V$ 上的线性变换。

任取非零向量 $\alpha\in V$，由于 $f(\varphi)(\alpha)=0$，根据\refpro{proposition:不可约多项式与向量组线性无关}知，向量组
\[
\{\alpha,\varphi(\alpha),\varphi^2(\alpha),\cdots,\varphi^{n-1}(\alpha)\}
\]
线性无关。而该向量组正是
\[
\{\varphi_1(\alpha),\varphi_2(\alpha),\cdots,\varphi_n(\alpha)\}.
\]\end{proof}











\end{document}